\section{Replacing an admissible theory by a prescribed finite skeleton}
\label{sec:136-finite-skeleton-replacement}

\paragraph{Archive and purpose.}
The requirements are cote136, physical pp.~117 and119,
footers116 and118, scans117 and119; common sources and
homotopy equalization are physical p.~121/footer120/scan121.
The following isolates the replacement part of a possible
common-second-skeleton construction. It does not construct
such a skeleton, assert that every prescribed skeleton qualifies,
or resolve the uniform integral common-source problem.
The algebraic construction uses strict limits and the elementary
normalized-chain description of a finite coskeleton. The latter
is the finite form of the Dold--Kan correspondence already used
in AQ-136-FIRST-SKELETON and ARCH-136-PATHS; its needed formula
and extension test are checked below.

Let \(r\geq0\) be finite. Let \(F\) be an \(r\)-truncated
simplicial, nonnegative bigraded strict differential PD
\(S=\Gamma_{\mathbb Z}(T)\)-algebra. Suppose
\(F_s^{0,*}=\mathbb Z[0]\) for \(0\leq s\leq r\),
with the specified coefficients, and assume the following tails.
For each form degree \(q\), there is a bound uniform over
\(0\leq s\leq r\) such that
\[
 H^0(F_s^{w,*})=S^w,\qquad
 H^q(F_s^{w,*})=0\quad(q>0)
                                                        \tag{136.FR1}
\]
in the respective weight ranges. The first equality means that
the coefficient map is the isomorphism. For each \(q\),
there is also a bound \(\mu_q\) such that the truncated
normalized row satisfies
\[
 H_i(NF_*^{w,q})=0\qquad(0\leq i<r,\ w\geq\mu_q).
                                                        \tag{136.FR2}
\]
When \(r=0\) the latter condition is empty. The homology
in (136.FR2) uses only the available chain terms through degree
\(r\). Let \(C\) be an admissible theory and let
\(u:F\to\operatorname{sk}_r C\) be a strict map preserving
coefficients and all the displayed structure.

\paragraph{ARCH-136-FINITE-SKELETON-REPLACEMENT.}
The strict pullback
\[
 C_F=C\mathbin{\times}_{\operatorname{cosk}_r\operatorname{sk}_r C}
                         \operatorname{cosk}_r F          \tag{136.FR3}
\]
is admissible, and its projection \(C_F\to C\) is a
strict arrow. Its truncation is literally \(F\), with that
projection restricting to \(u\). For each form degree, the
vertical and horizontal weight bounds are independent of the
simplex dimension. In particular, if one \(F\) satisfying
(136.FR1)--(136.FR2) maps to a finite family of admissible
\(C_j\), (136.FR3) gives admissible replacements of that
family with exactly the same \(r\)-skeleton.

\begin{proof}
Limits of strict differential PD algebras are computed in their
underlying bigraded groups, with products, differential and divided
powers componentwise. Thus both coskeleta and (136.FR3) are actual
strict objects. At a simplex of dimension at most \(r\), the
map \(C\to\operatorname{cosk}_r\operatorname{sk}_r C\)
is the identity after its canonical identification. The pullback
there is exactly \(F\). Nonnegative gradings and coefficient maps are retained. In weight
zero the map between the two coskeleta is an isomorphism, because
both truncated diagrams are the same coefficient diagram. The
pullback therefore has exactly the weight-zero part of \(C\).
This also handles \(r=0\), where a coskeleton itself has one
coefficient copy per vertex rather than only one copy.

We record the normalized-chain facts explicitly. For a truncated
simplicial module \(M\), set \(V=\operatorname{cosk}_r M\).
Its normalized complex is
\[
 (NV)_i=
 \begin{cases}
 (NM)_i,&0\leq i\leq r,\\
 Z_r(NM),&i=r+1,\\
 0,&i>r+1,
 \end{cases}
 \qquad d_{r+1}:Z_r(NM)\hookrightarrow(NM)_r.
                                                        \tag{136.FR4}
\]
For \(r=0\), the last map is the identity of \(M_0\).
Indeed a normalized matching family in dimension \(r+1\)
has every face but its zeroth face zero; the remaining face is
exactly a normalized \(r\)-cycle. The next matching relation
makes every higher normalized family zero. Together with the
usual direct decomposition into normalized and degenerate
summands, this proves (136.FR4). It follows that \(V\) is
acyclic exactly when the homology of \(NM\) below \(r\)
is zero.

The same decomposition gives the two surjectivity facts we need.
An untruncated simplicial module \(M\) maps onto
\(\operatorname{cosk}_r\operatorname{sk}_r M\) if
\(H_r(NM)=0\): the only new normalized surjectivity test is
\(d:N_{r+1}M\to Z_r(NM)\). For a truncated module, the
map from its \(s\)-simplices to compatible boundary values
on \(\partial\Delta[s]\), for \(1\leq s\leq r\),
is surjective when \(H_{s-1}(NM)=0\). Subtract the degeneracies
of all prescribed faces but one; the remaining obstruction is a
normalized \((s-1)\)-cycle, and a primitive supplies the missing
simplex. These are statements about modules and actual elements;
no chosen linear contraction is required.

Fix one form degree \(q\) and a sufficiently high weight for
(136.FR2) and the horizontal acyclicity of the corresponding
row of \(C\). Put
\(U=\operatorname{cosk}_r\operatorname{sk}_r C\) and
\(V=\operatorname{cosk}_r F\). Formula (136.FR4) makes
both rows \(U_*^{w,q}\) and \(V_*^{w,q}\) acyclic.
The map \(C_*^{w,q}\to U_*^{w,q}\) is surjective by
the preceding test. Consequently there is a short exact sequence
of simplicial modules
\[
 0\longrightarrow (C_F)_*^{w,q}
   \longrightarrow C_*^{w,q}\oplus V_*^{w,q}
   \xrightarrow{(a,b)\mapsto a-u(b)}U_*^{w,q}
   \longrightarrow0.                                  \tag{136.FR5}
\]
Normalization is exact, and all three other complexes are
acyclic. Its homology sequence proves the required horizontal
acyclicity of \(C_F\), including degree zero. The bound
has no simplex-dimension dependence.

Here is a uniform vertical comparison. Fix an integer
\(d\geq0\). Choose one weight bound with the following
properties: the maps \(F_s\to C_s\), for \(s\leq r\),
induce cohomology isomorphisms in degrees \(0,\ldots,d\);
all the row-boundary surjections just proved hold for both
truncated theories in form degrees \(0,\ldots,d+1\);
and \(C\to U\) is surjective in those form degrees.
Such a bound is a finite maximum of (136.FR1), (136.FR2),
and the original bounds of \(C\). In the first assertion,
the common coefficient identification in degree zero and the
vanishing of positive cohomology make every map the required
isomorphism. There is no condition on the number of simplices.

For every finite simplicial set \(K\) of dimension at most
\(r\), we claim that
\[
                         F(K)\longrightarrow C(K)
                                                        \tag{136.FR6}
\]
is a cohomology isomorphism through degree \(d\) in this
same weight range. Induct first on cell dimension, then on the number of cells in
that dimension. Build \(K\) by attaching its nondegenerate
simplices in that order. For a single attachment of
an \(s\)-simplex to \(K'\), sections give
\[
 0\longrightarrow F(K)
  \longrightarrow F(K')\oplus F_s
  \longrightarrow F(\partial\Delta[s])
  \longrightarrow0                                    \tag{136.FR7}
\]
in all form degrees \(0,\ldots,d+1\). The last map is
surjective because \(F_s\) has the boundary-extension
property proved above. There is the identical sequence for
\(C\). For zero-simplices the assertion is just the given
map \(F_0\to C_0\), with finite products for multiple
vertices. Apply the cohomology sequences after brutal truncation
in degree \(d+1\). These truncations retain cohomology through
\(d\), and the displayed degreewise surjections make them
short exact sequences of complexes. The induction hypothesis on
\(K'\) and on the boundary, together with the simplex comparison,
proves (136.FR6) by the five-term exact-sequence argument in each
degree through \(d\). This induction reuses the same bounds
at every cell; it does not enlarge them with the number of cells.

Take \(K=\operatorname{sk}_r\Delta[n]\). This is finite
for every \(n\), and the skeleton--coskeleton adjunction
identifies (136.FR6) with
\[
                                V_n\longrightarrow U_n.
                                                        \tag{136.FR8}
\]
Hence (136.FR8) is a vertical cohomology isomorphism through
\(d\) with one bound for all \(n\). In the guarded form
range, (136.FR5) is also a short exact sequence of vertical
complexes in each simplex, truncated in degree \(d+1\).
Compare it to the split exact sequence
\[
 0\longrightarrow C_n\longrightarrow C_n\oplus U_n
                 \longrightarrow U_n\longrightarrow0,
 \qquad x\longmapsto(x,x|_U),\quad (x,y)\longmapsto x|_U-y.
                                                        \tag{136.FR9}
\]
The middle map of the comparison is identity on \(C_n\)
and (136.FR8) on the other factor; the right map is identity.
The same five-term cohomology argument proves that
\((C_F)_n\to C_n\) is an isomorphism through degree
\(d\), uniformly in \(n\). The high-weight coefficient
cohomology and vertical vanishings of \(C\) now give those
of \(C_F\). Since \(d\) was arbitrary, this is exactly
the remaining admissibility condition. All maps preserve the
specified coefficient inclusion, so the degree-zero identification
is the required one, not a separate abstract isomorphism.
\end{proof}

\paragraph{What remains to construct.}
For \(r=2\), a common candidate must already have coefficient
vertical cohomology in all required weight tails in simplices
zero, one and two, and must fill the normalized row cycles in
degrees zero and one uniformly for each form degree. A partial
source with uncontrolled higher vertical cohomology in simplex
two does not meet (136.FR1). The theorem removes the subsequent
replacement step once these hypotheses are verified. It does
not supply that missing second skeleton or an iteration with
bounds uniform over all skeleton dimensions.

\paragraph{Verification status.}
The author completed a whole saved-proof reading, including
the case \(r=0\), the normalized matching formulas, all finite
row guards, brutal vertical truncation, the cell induction and
the pullback comparison. The cote15 peer then completed a separate whole saved-proof
reading and reported PASS, including the normalized formulas,
finite vertical truncations and uniform bounds. No common
second-skeleton construction is inferred.
